\subsection{Motivation}

The increasing accessibility of AI-based speech synthesis and voice cloning technologies has lowered the barrier for generating highly convincing synthetic speech. Modern voice generation systems are widely available through open-source frameworks, cloud services, and commercial applications, making it possible to create realistic voice impersonations with minimal technical expertise and limited reference audio.

The misuse of synthetic speech presents emerging challenges for cybersecurity, digital trust, and communication security. AI-generated voices can be used to impersonate individuals during phone calls, bypass voice-based verification mechanisms, conduct social engineering attacks, and spread misinformation through fabricated audio content. As voice communication continues to play an important role in online meetings, internet telephony, customer service systems, and voice authentication platforms, the potential impact of such attacks is increasing.

Although significant progress has been made in deepfake speech detection research, most existing solutions are evaluated using clean and carefully curated datasets that do not accurately represent real-world communication environments. Practical deployment requires detection systems that remain reliable under communication degradations and operational constraints encountered in everyday voice transmission systems. Therefore, there is a strong need to investigate detection approaches that can operate effectively in realistic communication scenarios while maintaining acceptable accuracy and latency.
